\section*{序論}
10月30日朝までの生成AI関連Xでは、フロンティアモデルの更新と、AIの行動原則を定める仕様の見直しが並んで話題になった。GoogleのGemini 4 Carbonはサイバーセキュリティと長時間タスクの性能を訴求し、OpenAIはModel Specを五つの中核的価値観を中心に再編した。モデルの能力だけでなく、利用時の制御や人間との役割分担も注目点となっている。

オープンウェイトと開発環境では、NVIDIAの多言語音声転写モデル、Appleデバイス向けローカル推論パッケージ、herdrのエージェント向けデスクトップ環境などが取り上げられた。GitHub Copilot関連の新機能やAnthropic APIの入力上限拡大も、開発者がAIを業務に組み込む際の選択肢を広げる話題として共有された。

本号はGrokが収集した12件を原文の順序で記事化する。Carbonの公式投稿は観測窓より前であり、期間内の議論と区別する。また、Metaの未発表モデルIDやエージェントのブラウザ脆弱性など、未確認の主張は確定事実として扱わない。独立した一次資料の再調査は実施していない。
